\documentclass{article}
\usepackage{amsmath,amssymb,amsthm}
\usepackage{algorithm}
\usepackage{algpseudocode}
\usepackage{hyperref}
\usepackage{geometry}
\geometry{margin=1in}
\newtheorem{theorem}{Theorem}
\newtheorem{lemma}{Lemma}
\newtheorem{assumption}{Assumption}
\newtheorem{corollary}{Corollary}
\newcommand{\Breg}[2]{D_{h}(#1,#2)}
\newcommand{\grad}{\nabla}
\newcommand{\prox}{\operatorname{prox}}
\begin{document}

\section*{Algorithm Name}
\textbf{Mirror Descent for Non‑Convex Smooth Functions (MD‑NC)}  

The method performs a Bregman‑proximal step using the gradient of the objective at the current iterate.

\section*{Mathematical Formulation}
Let $f:\mathbb{R}^{d}\rightarrow\mathbb{R}$ be differentiable and let $h:\mathbb{R}^{d}\rightarrow\mathbb{R}$ be a $\sigma$‑strongly convex \emph{prox‑function}.  
Denote the Bregman divergence induced by $h$ as  

\[
\Breg{x}{y}=h(x)-h(y)-\langle\nabla h(y),x-y\rangle .
\]

Given a stepsize sequence $\{\eta_k\}_{k\ge 0}$, the MD‑NC update is  

\[
\boxed{
x_{k+1}= \operatorname*{arg\,min}_{x\in\mathbb{R}^{d}}
\Big\{ \langle \grad f(x_k),\,x\rangle + \frac{1}{\eta_k}\,\Breg{x}{x_k}\Big\}
}
\tag{1}
\]

The optimality condition of (1) yields the \emph{gradient mapping}  

\[
G_{\eta_k}(x_k):=\frac{1}{\eta_k}\bigl(x_k-x_{k+1}\bigr)
          =\nabla h(x_k)-\nabla h(x_{k+1}) .
\tag{2}
\]

\section*{Pseudocode}
\begin{algorithm}[H]
\caption{Mirror Descent for Non‑Convex Smooth Functions}
\begin{algorithmic}[1]
\Require Initial point $x_0\in\mathbb{R}^{d}$, stepsize schedule $\{\eta_k\}_{k=0}^{T-1}>0$, prox‑function $h$ (strongly convex)
\For{$k=0$ \textbf{to} $T-1$}
    \State Compute stochastic (or full) gradient $g_k:=\grad f(x_k)$
    \State Solve the proximal subproblem  
    \[
    x_{k+1}= \arg\min_{x}\Big\{\langle g_k,x\rangle+\frac{1}{\eta_k}\Breg{x}{x_k}\Big\}
    \]
    \Comment{Closed‑form when $h$ is Euclidean: $x_{k+1}=x_k-\eta_k g_k$}
\EndFor
\State \textbf{output} $x_T$ (or the average $\frac1T\sum_{k=0}^{T-1}x_k$)
\end{algorithmic}
\end{algorithm}

\section*{Dry Run Example}
Consider the one‑dimensional quadratic  

\[
f(w)=(w-3)^2 ,\qquad \grad f(w)=2(w-3).
\]

Choose Euclidean prox $h(w)=\tfrac12 w^{2}$ (so $\Breg{u}{v}=\tfrac12 (u-v)^2$) and a constant stepsize $\eta_k=\eta=0.1$.  
Initialize $w_0=0$.

\begin{center}
\begin{tabular}{c|c|c|c}
Iteration $k$ & $g_k=\grad f(w_k)$ & $w_{k+1}=w_k-\eta g_k$ & $f(w_{k+1})$\\\hline
0 & $2(0-3)=-6$ & $0-0.1(-6)=0.6$ & $(0.6-3)^2=5.76$\\
1 & $2(0.6-3)=-4.8$ & $0.6-0.1(-4.8)=1.08$ & $(1.08-3)^2=3.6864$\\
2 & $2(1.08-3)=-3.84$ & $1.08-0.1(-3.84)=1.464$ & $(1.464-3)^2=2.3660$\\
\end{tabular}
\end{center}
The objective values decrease from $9$ (at $w_0$) to $2.37$ after three iterations, illustrating the descent property of (1).

\newpage
\section*{Assumptions}
\begin{assumption}[Strongly Convex Prox‑function]\label{as:h}
$h$ is $\sigma$‑strongly convex w.r.t. a norm $\|\cdot\|$:
\[
h(x)\ge h(y)+\langle\nabla h(y),x-y\rangle+\frac{\sigma}{2}\|x-y\|^{2},
\qquad\forall x,y\in\mathbb{R}^{d}.
\]
\end{assumption}

\begin{assumption}[Smoothness of $f$ w.r.t. $h$]\label{as:f}
$f$ is $L$‑smooth in the Bregman sense, i.e.
\[
f(x)\le f(y)+\langle\grad f(y),x-y\rangle+\frac{L}{2}\|x-y\|^{2},
\qquad\forall x,y\in\mathbb{R}^{d}.
\]
Equivalently,
\[
f(x_{k+1})\le f(x_k)+\langle\grad f(x_k),x_{k+1}-x_k\rangle
                +\frac{L}{2}\|x_{k+1}-x_k\|^{2}.
\tag{3}
\]
\end{assumption}

\begin{assumption}[Stepsize]\label{as:eta}
The stepsizes satisfy $0<\eta_k\le \frac{\sigma}{L}$ for all $k$.
\end{assumption}

\section*{Main Convergence Theorem}
\begin{theorem}[Non‑convex Mirror Descent]\label{thm:main}
Under Assumptions~\ref{as:h}--\ref{as:eta}, let $\{x_k\}_{k=0}^{T}$ be generated by \eqref{1}.  
Define the gradient mapping $G_{\eta_k}(x_k):=\frac1{\eta_k}(x_k-x_{k+1})$.  
Then the averaged squared norm satisfies  

\[
\frac{1}{T}\sum_{k=0}^{T-1}\|G_{\eta_k}(x_k)\|_{*}^{2}
\;\le\;
\frac{2L}{\sigma}\,\frac{f(x_0)-f^{\inf}}{T},
\tag{4}
\]
where $\|\cdot\|_{*}$ is the dual norm of $\|\cdot\|$ and $f^{\inf}:=\inf_{x}f(x)$.  
Consequently, there exists $k^{\star}\in\{0,\dots,T-1\}$ such that  

\[
\|G_{\eta_{k^{\star}}}(x_{k^{\star}})\|_{*}^{2}
\;\le\;
\frac{2L}{\sigma}\,\frac{f(x_0)-f^{\inf}}{T}.
\tag{5}
\]
Thus MD‑NC attains a sub‑linear rate $O(1/T)$ in terms of the gradient mapping norm.
\end{theorem}

\section*{Theoretical Properties}
\begin{itemize}
\item \textbf{Stability.} Because $h$ is strongly convex, the Bregman divergence controls the distance between iterates, guaranteeing that the sequence $\{x_k\}$ remains in a bounded level set of $f$.
\item \textbf{Hyperparameter Sensitivity.} The condition $\eta_k\le\sigma/L$ is essential: a larger stepsize would break the descent inequality (see Step~4 of the proof). In practice one may set $\eta=\sigma/L$ for the fastest guaranteed contraction.
\item \textbf{Comparison to Gradient Descent.} When $h(w)=\tfrac12\|w\|^{2}$, $\Breg{x}{y}=\tfrac12\|x-y\|^{2}$ and $G_{\eta_k}(x_k)=\grad f(x_k)$. The bound \eqref{4} reduces to the classical $O(1/T)$ guarantee for (deterministic) gradient descent on smooth non‑convex functions.
\item \textbf{Extension to Stochastic Gradients.} With an unbiased estimator $g_k$ of $\grad f(x_k)$ and bounded variance, the same analysis yields an $O(1/\sqrt{T})$ bound on the expected gradient mapping norm.
\end{itemize}

\section*{Discussion}
The theorem shows that even without convexity the mirror‑descent step provides a measurable progress certificate via the gradient mapping. The rate matches the best known for first‑order methods on smooth non‑convex problems, while the Bregman geometry allows the algorithm to exploit problem‑specific structure (e.g., simplex constraints with the negative entropy prox).

\newpage
\section*{Preliminaries (Definitions and Lemmas)}
\begin{lemma}[Three‑point identity for Bregman divergence]\label{lem:3pt}
For any $x,y,z\in\mathbb{R}^{d}$,
\[
\Breg{x}{z}= \Breg{x}{y}+ \Breg{y}{z}+ \langle \nabla h(y)-\nabla h(z),\,x-y\rangle .
\tag{6}
\]
\end{lemma}
\begin{proof}
Direct expansion of the definition of $\Breg{\cdot}{\cdot}$ yields \eqref{6}. \hfill $\square$
\end{proof}

\section*{Main Proof (Step‑by‑Step Derivation)}
We prove Theorem~\ref{thm:main}. Throughout we denote $g_k:=\grad f(x_k)$ and $\eta:=\eta_k$ (constant for brevity).

\begin{enumerate}
\item \textbf{Optimality condition of the subproblem.}  
From \eqref{1} and first‑order optimality,
\[
0\in g_k+\frac{1}{\eta}\bigl(\nabla h(x_{k+1})-\nabla h(x_k)\bigr)
\;\Longrightarrow\;
g_k = -\frac{1}{\eta}\bigl(\nabla h(x_{k+1})-\nabla h(x_k)\bigr).
\tag{Step 1}
\]
\item \textbf{Relate gradient mapping to Bregman divergence.}  
Using the definition \eqref{2} and the strong convexity of $h$,
\[
\|G_{\eta}(x_k)\|_{*}^{2}
= \frac{1}{\eta^{2}}\| \nabla h(x_k)-\nabla h(x_{k+1})\|_{*}^{2}
\ge \frac{\sigma^{2}}{\eta^{2}}\|x_{k+1}-x_k\|^{2}.
\tag{Step 2}
\]
The inequality follows from the $\sigma$‑strong convexity of $h$ which implies  
$\langle \nabla h(x_k)-\nabla h(x_{k+1}),x_k-x_{k+1}\rangle\ge\sigma\|x_k-x_{k+1}\|^{2}$ and the Cauchy–Schwarz inequality.

\item \textbf{Apply smoothness of $f$.}  
Using \eqref{3} with $y=x_k$ and $x=x_{k+1}$,
\[
f(x_{k+1})\le f(x_k)+\langle g_k, x_{k+1}-x_k\rangle+\frac{L}{2}\|x_{k+1}-x_k\|^{2}.
\tag{Step 3}
\]

\item \textbf{Replace $\langle g_k, x_{k+1}-x_k\rangle$ by the proximal optimality.}  
From Step~1,
\[
\langle g_k, x_{k+1}-x_k\rangle
= -\frac{1}{\eta}\langle \nabla h(x_{k+1})-\nabla h(x_k),\,x_{k+1}-x_k\rangle.
\tag{Step 4}
\]
By strong convexity of $h$,
\[
\langle \nabla h(x_{k+1})-\nabla h(x_k),\,x_{k+1}-x_k\rangle
\ge \sigma\|x_{k+1}-x_k\|^{2}.
\tag{Step 5}
\]
Insert \eqref{Step 4} and \eqref{Step 5} into \eqref{Step 3}:
\[
f(x_{k+1})\le f(x_k)-\frac{\sigma}{\eta}\|x_{k+1}-x_k\|^{2}
                +\frac{L}{2}\|x_{k+1}-x_k\|^{2}.
\tag{Step 6}
\]

\item \textbf{Choose stepsize satisfying $\eta\le\sigma/L$.}  
Because $\eta\le\sigma/L$, we have $\frac{L}{2}\le\frac{\sigma}{2\eta}$. Hence
\[
-\frac{\sigma}{\eta}+\frac{L}{2}
\le -\frac{\sigma}{2\eta}.
\tag{Step 7}
\]
Applying \eqref{Step 7} to \eqref{Step 6} yields the descent inequality
\[
f(x_{k+1})\le f(x_k)-\frac{\sigma}{2\eta}\|x_{k+1}-x_k\|^{2}.
\tag{Step 8}
\]

\item \textbf{Express the distance term via the gradient mapping.}  
From Step~2,
\[
\|x_{k+1}-x_k\|^{2}\le \frac{\eta^{2}}{\sigma^{2}}\|G_{\eta}(x_k)\|_{*}^{2}.
\tag{Step 9}
\]
Plugging \eqref{Step 9} into \eqref{Step 8} gives
\[
f(x_{k+1})\le f(x_k)-\frac{\eta}{2\sigma}\|G_{\eta}(x_k)\|_{*}^{2}.
\tag{Step 10}
\]

\item \textbf{Telescoping sum.}  
Summing \eqref{Step 10} from $k=0$ to $T-1$,
\[
f(x_T)\le f(x_0)-\frac{\eta}{2\sigma}\sum_{k=0}^{T-1}\|G_{\eta_k}(x_k)\|_{*}^{2}.
\tag{Step 11}
\]
Since $f(x_T)\ge f^{\inf}$, rearranging gives
\[
\frac{1}{T}\sum_{k=0}^{T-1}\|G_{\eta_k}(x_k)\|_{*}^{2}
\le \frac{2\sigma}{\eta T}\bigl(f(x_0)-f^{\inf}\bigr).
\tag{Step 12}
\]

\item \textbf{Insert the maximal admissible stepsize.}  
Set $\eta=\sigma/L$ (the largest stepsize allowed by Assumption~\ref{as:eta}). Then
\[
\frac{2\sigma}{\eta}= \frac{2\sigma}{\sigma/L}=2L.
\tag{Step 13}
\]
Substituting \eqref{Step 13} into \eqref{Step 12} yields exactly the bound \eqref{4} of Theorem~\ref{thm:main}.
\hfill $\square$
\end{enumerate}

\section*{Rate Derivation (With Constants)}
From \eqref{4},
\[
\min_{0\le k<T}\|G_{\eta_k}(x_k)\|_{*}^{2}
\le \frac{1}{T}\sum_{k=0}^{T-1}\|G_{\eta_k}(x_k)\|_{*}^{2}
\le \frac{2L}{\sigma}\,\frac{f(x_0)-f^{\inf}}{T}.
\]
Thus the algorithm achieves an $O(1/T)$ decay of the squared gradient‑mapping norm. When expressed in terms of the Euclidean norm (i.e., $h(w)=\tfrac12\|w\|^{2}$) this matches the standard non‑convex gradient‑descent rate.

\section*{Special Cases}
\begin{itemize}
\item \textbf{Euclidean setting.} $h(w)=\tfrac12\|w\|^{2}$, $\sigma=1$, $\Breg{x}{y}=\tfrac12\|x-y\|^{2}$. The update reduces to $x_{k+1}=x_k-\eta\grad f(x_k)$ and $G_{\eta}(x_k)=\grad f(x_k)$. Theorem~\ref{thm:main} recovers the classic result for gradient descent.
\item \textbf{Negative entropy on the simplex.} $h(x)=\sum_{i}x_i\log x_i$ (defined on the probability simplex). Here $\sigma=1$ w.r.t. the $\ell_1$ norm and the algorithm becomes the \emph{Exponentiated Gradient} method. The same $O(1/T)$ bound holds for the corresponding gradient mapping.
\item \textbf{Adaptive stepsizes.} If $\eta_k$ varies but always satisfies $\eta_k\le\sigma/L$, the same telescoping argument works with $\eta_k$ inside the sum, giving  
\[
\sum_{k=0}^{T-1}\eta_k\|G_{\eta_k}(x_k)\|_{*}^{2}
\le 2\sigma\bigl(f(x_0)-f^{\inf}\bigr).
\]
Choosing $\eta_k=\sigma/(L\sqrt{k+1})$ yields an $O(1/\sqrt{T})$ bound for the averaged mapping norm.
\end{itemize}

\end{document}